% Page 061 — page imprimée 58
% Georges Roth — suite, témoignage sur Olivier Giran

Dénoncé par deux traîtres qui avaient réussi à s’introduire dans le réseau, il fut arrêté par la
Gestapo le 30 juin 1942 en gare de Dijon. Emprisonné pendant de longs mois à il ne dénonça
personne et ne lâcha jamais le nom de Michel Hollard. Finalement il fut transféré à Fresnes
puis à Angers où il fut fusillé le 16 Avril 1943. Désormais une place d’Angers porte son nom.
Il avait pu écrire à sa famille un certain nombre de lettres qui ont été conservées. Voici le
dernier message d’adieu qu’il termina en son dernier matin, avant que les gardiens ne
viennent le chercher :

\begin{quote}\itshape
« ...Je suis calme, mais la mort me presse et j’ai tant à vous dire. D’abord, réglons bien certaines choses :
devant Dieu je me présente conscient et confiant. Je crois en lui, je crois en la vie de l’âme, en une vie spirituelle
où il me semble, où je suis sûr même, c’est mon malheureux cœur qui m’en assure, que je communierai avec
vous par la pensée dans tout le beau, le noble le juste auquel je vais me trouver mêlé.
Devant les hommes j’ai fait avec élan, joie, enthousiasme à certains moments, ce que j’ai cru de mon devoir de
faire. C’est la guerre : je tombe après d’autres, avant d’autres encore.

Je suis heureux de ma destinée, plus heureux que beaucoup qui végètent et s’ennuient... Gardez à jamais de moi
une image souriante, gaie, vivante d’un fils qui, grâce à vous et avec la con science de l’amour profond qui nous
unit, meurt heureux...

Oui ! Vive la France ! Les hommes sont lâches, traîtres ou salauds. La France est pure, propre et vivante, je suis
heureux, je ne meurs ni pour une faction, ni pour un homme, mais pour mon idée à moi de la servir, pour elle
pour mon patrimoine, pour vous deux que j’adore.
Je suis heureux, je vous aime

\begin{flushright}
On ouvre\\
Adieu ! Je vous embrasse fort\\
Olivier
\end{flushright}
\end{quote}

L’aumônier allemand qui a accompagné Olivier en ses derniers instants a écrit à sa mère, Mme.
Giran, une longue lettre :

Très honorée Mme. Giran,
...Mes notes sur la mort de votre fils et les dernières heures de sa vie ont été perdues lorsque j’ai
été fait prisonnier. Mais l’impression profonde que m’a laissée ma rencontre avec cet être d’élite
ne s’est pas estompée et je peux ainsi, aujourd’hui encore, vous envoyer une image vivante de ce
qu’ont été ses dernières heures…
Je lui dis que je regrettais qu’il me soit défendu de mettre à sa disposition un pasteur français et je
le priai de voir en moi le pasteur et non l’allemand : J’ai seulement le désir sincère de pouvoir
vous aider en tant que pasteur… Il me remercia et dit :

\begin{quote}\itshape
« Je dois dire que je ne considère pas la mort que je dois subir comme un malfaiteur. Je meurs pour la
France et pour mon peuple... Je savais ce que je faisais. Je n’ai pas trahi mais j’ai, au contraire, combattu
avec toutes mes forces contre les ennemis de mon pays et pour la liberté de la France... »
\end{quote}

Il parlait sans emphase. Mais la profonde sincérité et le grand amour de son pays qui se
traduisaient par ces mots me firent une profonde impression…. Puis je me mis à prier pour lui et
pour ses parents également. Il dit le « Notre Père » à haute voix. Les derniers mots de la prière
furent prononcés par lui avec la plus profonde émotion… Tout à coup il se tourna vers moi, saisit
ma main et me dit :

\begin{quote}\itshape
« Voyez-vous, toute cette éclatante éternité m’appartient. Le pas que je fais dans la mort n’est qu’un pas
dans cette immensité, dans la vraie liberté. Ce que je perds est peu de chose comparé à ce que j’acquiers
par ma mort. »
\end{quote}
